\section{Introducción al Dossier de Ingeniería de Hardware}

Este dossier presenta el análisis detallado del hardware del collar de geocercamiento DR-VF-01, incluyendo la arquitectura del sistema, el presupuesto de potencia, el análisis térmico y la lista de materiales (BOM).

\section{Arquitectura del Sistema}

La arquitectura del sistema se basa en el microcontrolador ESP32-S3-WROOM-1U-N16R8, que integra un procesador dual-core Xtensa LX7 con 16 MB de memoria Flash y 8 MB de PSRAM. Esta plataforma proporciona la capacidad de procesamiento necesaria para ejecutar algoritmos de geocercamiento en tiempo real y gestionar la pila de comunicaciones LoRa.

\subsection{Componentes Principales}

\begin{itemize}[leftmargin=*]
    \item \textbf{Microcontrolador Principal}: ESP32-S3-WROOM-1U-N16R8 (16 MB Flash, 8 MB PSRAM).
    \item \textbf{Sistema de Posicionamiento}: u-blox MAX-M10S GNSS multiconstelación.
    \item \textbf{Sensor de Inercia y Movimiento}: LSM6DSOX IMU de 6 ejes con bajo consumo.
    \item \textbf{Transceptor de Radiofrecuencia}: SX1262 LoRa Sub-GHz (915 MHz).
    \item \textbf{Gestión de Energía y Carga}: Cargador solar MPPT CN3791 (6 V) + LDO ultra-bajo \textit{Iq} TPS7A0533 (3.3 V).
\end{itemize}

\section{Presupuesto de Potencia}

\subsection{Metodología de Cálculo}

El presupuesto de potencia se calculó considerando el perfil de ciclo de trabajo (\textit{duty cycle}) en los distintos estados operativos del collar:

\begin{enumerate}[leftmargin=*]
    \item \textbf{Estado de Sueño Profundo (\textit{Deep Sleep})}: 98.5\% del tiempo operativo.
    \item \textbf{Verificación de Geocerca (GNSS Fix)}: 10 iteraciones por día.
    \item \textbf{Transmisión LoRa (Telemetría)}: 20 ráfagas por día.
    \item \textbf{Activación de Estímulo Disuasivo}: 2 eventos por día.
\end{enumerate}

\subsection{Resultados del Análisis de Consumo}

\begin{table}[H]
\centering
\caption{Desglose del consumo de corriente en modo \textit{Deep Sleep}}
\label{tab:power_deepsleep}
\vspace{2mm}
\begin{tabularx}{\textwidth}{X l r r}
\toprule
\rowcolor{categorybg}
\textbf{Componente} & \textbf{Estado Operativo} & \textbf{Corriente Base} & \textbf{Contribución Promedio} \\
\midrule
ESP32-S3 & Deep Sleep & 5.0 $\mu$A & 4.93 $\mu$A \\
\rowcolor{rowalt}
u-blox MAX-M10S & Standby & 25.0 $\mu$A & 24.63 $\mu$A \\
LSM6DSOX IMU & Low-Power & 10.0 $\mu$A & 9.85 $\mu$A \\
\rowcolor{rowalt}
SX1262 LoRa & Sleep & 0.7 $\mu$A & 0.69 $\mu$A \\
CN3791 MPPT & Standby & 150.0 $\mu$A & 150.00 $\mu$A \\
\bottomrule
\end{tabularx}
\end{table}

\subsection{Autonomía de Batería}

Con una celda LiPo de 1500 mAh y una corriente promedio efectiva de 0.362 mA, la autonomía teórica estimada del dispositivo se calcula como:

$$\text{Autonomía (días)} = \frac{1500\text{ mAh}}{0.362\text{ mA} \times 24\text{ h/día}} \approx 173\text{ días}$$

Considerando la autodescarga de la batería y factores de degradación ambiental, se estima una operación continua superior a 147 días en condiciones ideales (capacidad efectiva de 1275 mAh, con factor de degradación de 0.85).

\section{Análisis Térmico}

\subsection{Componentes de Alta Disipación}

\begin{itemize}[leftmargin=*]
    \item \textbf{Regulador LDO TPS7A0533}: 0.255 W disipados en regulación pico.
    \item \textbf{Cargador Solar CN3791}: 0.140 W disipados en carga máxima.
    \item \textbf{Transceptor LoRa SX1262}: 0.396 W durante eventos de transmisión en potencia máxima (+22 dBm).
\end{itemize}

\subsection{Cálculos de Temperatura de Juntura}

Para el regulador LDO TPS7A0533 con encapsulado SOT-23 ($R_{\theta\text{JA}} = 60.9\,^\circ\text{C/W}$):

$$P_{\text{disipada}} = (V_{\text{in}} - V_{\text{out}}) \times I_{\text{out}} = (5.0\text{ V} - 3.3\text{ V}) \times 0.15\text{ A} = 0.255\text{ W}$$

$$\Delta T = P_{\text{disipada}} \times R_{\theta\text{JA}} = 0.255\text{ W} \times 60.9\,^\circ\text{C/W} = 15.53\,^\circ\text{C}$$

$$T_{\text{juntura}} = T_{\text{ambiente}} + \Delta T = 65.0\,^\circ\text{C} + 15.53\,^\circ\text{C} = 80.53\,^\circ\text{C}$$

Dado que la temperatura máxima de juntura permitida es de $125\,^\circ\text{C}$, el componente opera dentro del margen de seguridad del 35\%.

\section{Lista de Materiales (BOM)}

\subsection{Componentes Electrónicos Principales}

\begin{table}[H]
\centering
\caption{Lista de componentes principales (BOM) del hardware DR-VF-01}
\label{tab:hardware_bom}
\vspace{2mm}
\begin{tabularx}{\textwidth}{c X l l l}
\toprule
\rowcolor{categorybg}
\textbf{Cant.} & \textbf{Descripción} & \textbf{MPN} & \textbf{Fabricante} & \textbf{Proveedor} \\
\midrule
1 & ESP32-S3 (16 MB Flash, 8 MB PSRAM) & ESP32-S3-WROOM-1U-N16R8 & Espressif & Digi-Key \\
\rowcolor{rowalt}
1 & Módulo GNSS Multiconstelación & MAX-M10S-00B & u-blox & Avnet \\
1 & Sensor de Inercia IMU 6 ejes & LSM6DSOXTR & STMicroelectronics & Digi-Key \\
\rowcolor{rowalt}
1 & Transceptor LoRa Sub-GHz & SX1262IMLTRT & Semtech & Digi-Key \\
\bottomrule
\end{tabularx}
\end{table}

\subsection{Estructura de Costos Estimada}

El costo total directo por unidad producida (lote de 100 unidades) es de \$43.25 USD:

\begin{itemize}[leftmargin=*]
    \item \textbf{Componentes Electrónicos (PCBA)}: \$24.50 USD
    \item \textbf{Componentes Mecánicos y Carcasa IP67}: \$8.75 USD
    \item \textbf{Fabricación de PCB (4 capas, FR4)}: \$3.25 USD
    \item \textbf{Ensamblaje y SMT}: \$5.50 USD
    \item \textbf{Empaque y Control de Calidad}: \$1.25 USD
\end{itemize}

\section{Conclusiones}

El diseño de hardware del collar DR-VF-01 se orienta a cumplir los requerimientos operativos y de autonomía fijados; sin embargo, el consumo calculado en \emph{Deep Sleep} ($190\,\mu\text{A}$) supera el límite de $150\,\mu\text{A}$ de \texttt{docs/REQUIREMENTS.md} (RNF-03) y está pendiente de revisión. La selección de un regulador LDO de ultra-bajo consumo en combinación con el microcontrolador ESP32-S3 garantiza una operación extendida sin comprometer la velocidad de procesamiento requerida para las geocercas en tiempo real.